\chapter*{Resumen}

\section*{\tituloPortadaVal}

El transformador (\textit{transformer}) es una aquitectura utilizada en las redes neurales que
apareció en 2017 y ha revolucionado el campo del procesamiento de lenguaje natural (con los
modelos extensos de lenguaje), además de generar resultados excelentes en otras áreas como
la visión por computadora o en tareas que requieren la combinación de múltiples tipos de datos
(imágenes, texto, audio, vídeo). Se basan en un concepto de procesamiento conocido como
“atención” que permite a la red asignar diferentes pesos a entradas distintas.


En este trabajo se pretende estudiar esta arquitectura y presentarla de manera didáctica,
analizando cuáles son las intuiciones que la motivan y las justificaciones matemáticas, hasta
donde las haya, que la hacen posible. También se realizarán implementaciones cuya naturaleza
puede variar conforme se vaya realizando el trabajo, pero que en principio estarían orientadas a
ilustrar cómo cambios en la arquitectura y/o en el valor de los pesos pueden afectar el
rendimiento del modelo.



   


